\section{Introduction}
LLM‑based agents have rapidly moved from controlled demonstrations to real‑world deployments, supporting complex multi‑step tasks in domains such as autonomous software engineering \citep{wu2023autogen}, open‑world tool use \citep{shen2024hugginggpt}, and everyday personal assistance \citep{ge2024openagi}. Although LLM‑based agents have achieved widespread adoption, practitioners consistently report substantial end‑to‑end latency \citep{hua2024interactive, zhang2024ecoact, zhang2023ecoassistant, saha2024system}. This latency not only degrades user experience but also limits the feasibility of deploying LLM‑based agents in time‑sensitive scenarios such as real‑time decision support \citep{cai2025rtbagent} or interactive tutoring \citep{modran2024llm, piro2024mylearningtalk}.

Multiple lines of research attempt to address the latency bottlenecks. These include approaches such as context length reduction \citep{zhang2024ecoact}, step parallelization \citep{hua2024interactive, li2025parallelized, pan2025learning}, and dual-process thinking that combines fast and deliberate reasoning modes \citep{saha2024system, lin2023swiftsage, christakopoulou2024agents}. While these approaches offer various improvements in efficiency, there are several common critical limitations: (1) \textbf{No guarantee on lossless performance} except ISP: \cite{saha2024system} requires training a router between two agents and \cite{zhang2024ecoact} requires careful prompt engineering to enable runtime tool-selection, both relying on the quality of the additional component. \cite{pan2025learning} requires the model to learn to generate parallelizable steps. (2) \textbf{Pre-deployment preparation}: \cite{saha2024system, pan2025learning} require training specialized modules in the system before deployment, adding complexity and overhead to implementation. (3) \textbf{Lack of user-controllability} \citep{kiefer2017controllability}: Most existing approaches neglect user-controlled flexibility in navigating efficiency improvements. This inflexibility becomes particularly problematic in the rapidly evolving LLM ecosystem, where pricing structures, model inference speeds, and model performance change frequently, and organizations have diverse priorities regarding latency versus operational costs including router training, prompt engineering, or extra cost brought by parallelizing computations.


In this paper, we propose Dynamic Speculative Planning (DSP), a lossless, user-controllable agent planning acceleration framework that requires no pre-deployment preparation. In addition to ensuring lossless agentic performance via speculative strategy \citep{hua2024interactive}, DSP introduces three core improvements for cost reduction and user-controllability:

First, we \textbf{investigate the Pareto frontier} of the latency-cost trade-off in speculative agent planning, determining the optimal cost for any given latency reduction target. This trade-off is governed by the \emph{speculation step}, which refers to the number of future steps that an approximation agent attempts to predict ahead of verification. It determines how aggressively the system parallelizes computation to reduce latency. Our analysis reveals that the commonly adopted fixed-speculation-step approach \citep{hua2024interactive, leviathan2023fast} is fundamentally limited: for complex tasks, aggressive speculation produces excessive and redundant agent calls that substantially increase costs, while for simpler tasks, conservative speculation fails to deliver sufficient acceleration. Since optimal speculation step is highly context-dependent and unknown a priori, fixed steps significantly constrain practical utility.

Second, we develop \textbf{a lightweight adaptive speculation step predictor} that dynamically determines when to suspend speculation, effectively eliminating unnecessary cost while preserving acceleration benefits. This predictor employs \textit{online reinforcement learning} that requires no external datasets, preprocessing steps, or dedicated pre-deployment phases. Instead, the system learns to optimize speculation step organically as it processes tasks, becoming increasingly efficient over time with zero additional infrastructure costs while ensuring immediate deployment benefits.

Third, based on this dynamic prediction module, we introduce \textbf{two mechanisms that modulate the trade-off between latency and cost}: (1) biased step prediction and (2) unbiased step with biased offset, allowing practitioners to precisely calibrate system behavior anywhere along the spectrum from low-cost/low-speed to high-cost/high-speed operation. By providing fine-grained user control over the latency-cost balance, our framework accommodates diverse organizational priorities and adapts to the fluctuating pricing structures and inference speeds of the rapidly evolving LLM ecosystem.

\begin{figure}
    \centering
    % \small
    \resizebox{0.762\linewidth}{!}{
        \begin{tikzpicture}
         % Add legend at the bottom of the graph (one line, full width)
    \node[anchor=north, draw, fill=white, text width=12cm, inner sep=5pt, font=\scriptsize] 
    at (-0.7,-4) {
    \tikz\draw[fill=lightpurple!70, draw=lightpurple!70] (0,0) rectangle (0.3,0.15); \textcolor{purple!80!black}{Predictor} \hfill
    \tikz\draw[fill=lightblue!50, draw=lightblue!50] (0,0) rectangle (0.3,0.15); \textcolor{blue!80!black}{Approximation} \hfill
    \tikz\draw[fill=lightorange!50, draw=lightorange!50] (0,0) rectangle (0.3,0.15); \textcolor{orange!80!black}{Target} \hfill
    \tikz\draw[fill=almond!50, draw=almond!50] (0,0) rectangle (0.3,0.15); \textcolor{almond!80!black}{Trainer} \hfill
    \tikz\draw[fill=black!30, draw=black!30] (0,0) rectangle (0.3,0.15); Action execution \hfill
    \tikz\draw[pattern=north west lines, pattern color=lightblue!100] (0,0) rectangle (0.3,0.15); Canceled operations
};
    
        % Add this block for the Speculation Depth Predictor
        \node (pred1) at (-5,1.0) [rectangle, draw=lightpurple!70, fill=lightpurple!70, minimum width=10pt, minimum height=10pt, inner sep=0pt, anchor=west] {};
        \node at (-6.7,1.0) {\textcolor{purple!80!black}{Predictor}};
        
        % Draw a connecting line or arrow
        %\draw [->, dashed, purple!80!black] (pred.south) -- ++(0,-0.8);

        \node (plan1) at (pred1.east) [xshift=15pt] {\textcolor{purple!80!black}{\footnotesize \makecell{$k=2$}}};

        % Add this block for the second Speculation Depth Predictor
        \node (pred2) at (-1.98,1.0) [rectangle, draw=lightpurple!70, fill=lightpurple!70, minimum width=10pt, minimum height=10pt, inner sep=0pt, anchor=west] {};

        \node (plan2) at (pred2.east) [xshift=15pt] {\textcolor{purple!80!black}{\footnotesize \makecell{$k=3$}}};
        
        % approximation code
        \node (step1) at (-5,0) [rectangle, draw=lightblue!50, fill=lightblue!50, minimum width=25pt, minimum height=12pt, inner sep=0pt, anchor=west] {};
    
        \node (step1) at (-5,0) [rectangle, draw=lightblue!50, fill=lightblue!50, minimum width=25pt, minimum height=12pt, inner sep=0pt, anchor=west] {};
        \node (exc1) at (step1.east) [rectangle, draw=black!30, fill=black!30, minimum width=5pt, minimum height=12pt, inner sep=0pt, anchor=west] {};
        \node (step2) at (exc1.east) [rectangle, draw=lightblue!50, fill=lightblue!50, minimum width=20pt, minimum height=12pt, inner sep=0pt, anchor=west] {};
        \node (exc2) at (step2.east) [rectangle, draw=black!30, fill=black!30, minimum width=5pt, minimum height=12pt, inner sep=0pt, anchor=west] {};
        \node at (-6.3,0) {\textcolor{blue!80!black}{Approximation}};

        \node (plan1) at (step1.east) [yshift=15pt] {\textcolor{blue!80!black}{\scriptsize \makecell{A\\$\downarrow$}}};
        \node (plan2) at (step2.east) [yshift=15pt] {\textcolor{blue!80!black}{\scriptsize \makecell{B\\$\downarrow$}}};

        % target code
        \node (tar_step1) at ([yshift=-30pt]step1.west) [rectangle, draw=lightorange!50, fill=lightorange!50, minimum width=58pt, minimum height=12pt, inner sep=0pt, anchor=west] {};
        \node (tar_step2) at ([yshift=-45pt]step2.west) [rectangle, draw=lightorange!50, fill=lightorange!50, minimum width=55pt, minimum height=12pt, inner sep=0pt, anchor=west] {};
        \node at (-6.9,-1.6) {\textcolor{orange!80!black}{Target}};

        \draw [-, dashed, thick] ([yshift=7pt]pred1.west) -- node [yshift=-59.5pt] {\footnotesize $s_1$} ++(0,-3.75);
        \draw [-, dashed, thick] ([yshift=10pt]step2.west) -- node [yshift=-46pt] {\footnotesize $s_2$} ++(0,-2.9);
        \draw [-, dashed, thick] ([yshift=7pt]pred2.west) -- node [yshift=-59pt] {\footnotesize $s_3$} ++(0,-3.75);

        \node (tar_plan1) at (tar_step1.east) [yshift=15pt] {\textcolor{orange!80!black}{\scriptsize \makecell{A \textcolor{green!80!black}{\checkmark}\\$\downarrow$}}};

        % approximation code for episode 2
        \node (step3) at (tar_step2.east) [yshift=45pt] [rectangle, draw=lightblue!50, fill=lightblue!50, minimum width=25pt, minimum height=12pt, inner sep=0pt, anchor=west] {};
        \node (exc3) at (step3.east) [rectangle, draw=black!30, fill=black!30, minimum width=5pt, minimum height=12pt, inner sep=0pt, anchor=west] {};
        \node (step4) at (exc3.east) [rectangle, draw=lightblue!50, pattern=north west lines, pattern color=lightblue!100, minimum width=30pt, minimum height=12pt, inner sep=0pt, anchor=west] {};
        \node (exc4) at (step4.east) [rectangle, draw=black!30, pattern=north west lines, pattern color=black!50, minimum width=5pt, minimum height=12pt, inner sep=0pt, anchor=west] {};
        \node (step5) at (exc4.east) [rectangle, draw=lightblue!50, pattern=north west lines, pattern color=lightblue!100, minimum width=11pt, minimum height=12pt, inner sep=0pt, anchor=west] {};
        
        \node (plan3) at (step3.east) [yshift=15pt] {\textcolor{blue!80!black}{\scriptsize \makecell{C'\\$\downarrow$}}};
        \node (plan4) at (step4.east) [yshift=15pt] {\textcolor{blue!80!black}{\scriptsize \makecell{D'\\$\downarrow$}}};
        % \node (plan5) at (step5.east) [yshift=15pt] {\textcolor{blue!80!black}{\scriptsize \makecell{E\\$\downarrow$}}};

        % target code
        \node (tar_step3) at ([yshift=-30pt]step3.west) [rectangle, draw=lightorange!50, fill=lightorange!50, minimum width=78pt, minimum height=12pt, inner sep=0pt, anchor=west] {};
        \node (tar_exc3) at (tar_step3.east) [rectangle, draw=black!30, fill=black!30, minimum width=5pt, minimum height=12pt, inner sep=0pt, anchor=west] {};
        \node (tar_plan2) at (tar_step2.east) [yshift=15pt, xshift=-10pt] {\textcolor{orange!80!black}{\scriptsize \makecell{B \textcolor{green!80!black}{\checkmark}\\$\downarrow$}}};
        \node (tar_step4) at ([yshift=-45pt]step4.west) [rectangle, draw=lightorange!50, pattern=north west lines, pattern color=lightorange!100, minimum width=48pt, minimum height=12pt, inner sep=0pt, anchor=west] {};
        \node (tar_step5) at ([yshift=-60pt]step5.west) [rectangle, draw=lightorange!50, pattern=north west lines, pattern color=lightorange!100, minimum width=11pt, minimum height=12pt, inner sep=0pt, anchor=west] {};

        \draw [->, thick] (-7,-2.5) -- node [below, xshift=160pt, yshift=-3pt] {Time} (5.5,-2.5);
        \draw [-, dashed, thick] ([yshift=8pt]step4.west) -- node [yshift=-46pt] {} ++(0,-2.75);
        \draw [-, dashed, thick] ([yshift=8pt]step5.west) -- node [yshift=-46pt] {} ++(0,-2.77);

        \node (tar_plan3) at (tar_step3.east) [yshift=15pt, xshift=-22pt] {\textcolor{orange!80!black}{\scriptsize \makecell{C \textcolor{red}{\ding{55}}\\$\downarrow$}}};

        % online learning part
        \node (online) at ([yshift=75pt]tar_exc3.east) [rectangle, draw=almond!50, fill=almond!50, minimum width=78pt, minimum height=12pt, inner sep=0pt, anchor=west] {};
        \node (online_action) at (online.east) [yshift=18pt, xshift=-30pt] {\textcolor{orange!80!black}{\scriptsize \makecell{New data for Training:}}};
        \node (online_data) at (online.east) [yshift=10pt, xshift=27pt] {\textcolor{orange!80!black}{\scriptsize \makecell{($s = s_1, k = 3$)\\($s = s_2, k = 2$)\\($s = s_3, k = 1$)}}};

        \node at (-6.85,1.65) {\textcolor{almond!80!black}{Trainer}};

        % k = 3
        \node (pred3) at (tar_exc3.east) [yshift=58pt, rectangle, draw=lightpurple!70, fill=lightpurple!70, minimum width=10pt, minimum height=10pt, inner sep=0pt, anchor=west] {};
        \node (plan3) at (pred3.east) [xshift=15pt] {\textcolor{purple!80!black}{\footnotesize \makecell{$k=2$}}};

        % approximation code for episode 3
        \node (step6) at (tar_exc3.east)  [yshift=30pt, rectangle, draw=lightblue!50, fill=lightblue!50, minimum width=25pt, minimum height=12pt, inner sep=0pt, anchor=west] {};
        % cancel words
        \draw [-, solid, thick, draw=red, line width=1.5pt] ([yshift=6pt]step5.east) -- node [yshift=-55pt] {\footnotesize \textcolor{red}{\makecell{Cancel ongoing threads}}} ++(0,-2.7);
        % restart episode 3
        \node (exc6) at (step6.east) [rectangle, draw=black!30, fill=black!30, minimum width=5pt, minimum height=12pt, inner sep=0pt, anchor=west] {};
        \node (step7) at (exc6.east) [rectangle, draw=lightblue!50, fill=lightblue!50, minimum width=30pt, minimum height=12pt, inner sep=0pt, anchor=west] {};
        \node (exc7) at (step7.east) [rectangle, draw=black!30, fill=black!30, minimum width=5pt, minimum height=12pt, inner sep=0pt, anchor=west] {};
        

        \node (plan6) at (step6.east) [yshift=15pt] {\textcolor{blue!80!black}{\scriptsize \makecell{D\\$\downarrow$}}};
        \node (plan7) at (step7.east) [yshift=15pt] {\textcolor{blue!80!black}{\scriptsize \makecell{E\\$\downarrow$}}};

        % target code for episode 3
        \node (tar_step6) at (tar_exc3.east) [rectangle, draw=lightorange!50, fill=lightorange!50, minimum width=75pt, minimum height=12pt, inner sep=0pt, anchor=west] {};
        \node (tar_step7) at (exc6.east) [yshift=-45pt, rectangle, draw=lightorange!50, fill=lightorange!50, minimum width=63pt, minimum height=12pt, inner sep=0pt, anchor=west] {};

        \node (tar_plan6) at (tar_step6.east) [yshift=15pt] {\textcolor{orange!80!black}{\scriptsize \makecell{D \textcolor{green!80!black}{\checkmark}\\$\downarrow$}}};
        \node (tar_plan7) at (tar_step7.east) [yshift=15pt] {\textcolor{orange!80!black}{\scriptsize \makecell{E \textcolor{green!80!black}{\checkmark}\\$\downarrow$}}};

        \draw [-, dashed, thick] ([yshift=8pt]step5.west) -- node [yshift=-46pt] {} ++(0,-2.77);

        \draw [-, dashed, thick] ([yshift=8pt]step7.west) -- node [yshift=-45pt] {\footnotesize $s_5$} ++(0,-2.8);

        \draw [<->, thick]([yshift=5pt]step1.west |- 0,-3.5) -- node [below] {Total task time}   ([yshift=5pt]tar_step7.east |- 0,-3.5);

        \draw [-, dashed, thick] ([yshift=9pt]online.west) -- node [yshift=-69pt] {\footnotesize $s_4$} ++(0,-4.4);


        % curly bracket
        \draw[decoration={brace,raise=75pt},decorate] ([xshift=-2pt]step1.west) -- node[above=75pt] {episode 1} ([xshift=-5pt]step5.east);

        % \draw[decoration={brace,raise=105pt},decorate] (tar_step3.west) -- node[above=105pt] {episode 2} ([xshift=-5pt] tar_exc3.east);

        \draw[decoration={brace,raise=30pt},decorate] (online.west) -- node[above=30pt] {episode 2} ([yshift=90pt, xshift=2pt]tar_step7.east);
        
    \end{tikzpicture}
    }
    \caption{Dynamic Speculative Planning with Online Reinforcement Learning. The diagram illustrates how DSP adaptively adjusts speculation steps across planning episodes. (1) \textbf{Predictor} dynamically infers step values $k$ from current state $s_i$, while the \textbf{Approximation $(\approxi)$} and \textbf{Target $(\target)$} agents execute parallelly. (2) When mismatches occur (denoted by \textcolor{red}{\ding{55}}), ongoing threads are canceled and execution resumes from the last verified step. (3) An asynchronous \textbf{Trainer} collects state-step ($s, k$) pairs from completed episodes, continuously updating the predictor without blocking execution.}
    \label{fig:main_graph}
    % \vspace{-20pt}
    \vspace{-4.5pt}
    
    % \setlength{\belowcaptionskip}{-4pt} 
\end{figure}

\setlength{\textfloatsep}{3pt plus 2pt minus 2pt}

The structure of our system is presented in Figure \ref{fig:main_graph}. Our empirical evaluation across two standard agent benchmarks demonstrate that DSP achieves the efficiency of the fastest lossless acceleration methods, while reducing total cost by up to 30\%, cutting unnecessary cost by 60\%, and preserving comparable acceleration, demonstrating the effectiveness of the framework. 

%The contributions of this paper are threefold: (1) a systematic analysis of cost inefficiencies in speculative planning for LLM agents; (2) a novel dynamic speculation approach that approaches the theoretical Pareto frontier of the time-cost tradeoff; and (3) the first user-controllable acceleration framework for LLM-based agents that accommodates varying deployment priorities and economic constraints. 

% \begin{figure}[t]
% \centering
% \begin{figure}
    \centering
    % \small
    \resizebox{0.762\linewidth}{!}{
        \begin{tikzpicture}
         % Add legend at the bottom of the graph (one line, full width)
    \node[anchor=north, draw, fill=white, text width=12cm, inner sep=5pt, font=\scriptsize] 
    at (-0.7,-4) {
    \tikz\draw[fill=lightpurple!70, draw=lightpurple!70] (0,0) rectangle (0.3,0.15); \textcolor{purple!80!black}{Predictor} \hfill
    \tikz\draw[fill=lightblue!50, draw=lightblue!50] (0,0) rectangle (0.3,0.15); \textcolor{blue!80!black}{Approximation} \hfill
    \tikz\draw[fill=lightorange!50, draw=lightorange!50] (0,0) rectangle (0.3,0.15); \textcolor{orange!80!black}{Target} \hfill
    \tikz\draw[fill=almond!50, draw=almond!50] (0,0) rectangle (0.3,0.15); \textcolor{almond!80!black}{Trainer} \hfill
    \tikz\draw[fill=black!30, draw=black!30] (0,0) rectangle (0.3,0.15); Action execution \hfill
    \tikz\draw[pattern=north west lines, pattern color=lightblue!100] (0,0) rectangle (0.3,0.15); Canceled operations
};
    
        % Add this block for the Speculation Depth Predictor
        \node (pred1) at (-5,1.0) [rectangle, draw=lightpurple!70, fill=lightpurple!70, minimum width=10pt, minimum height=10pt, inner sep=0pt, anchor=west] {};
        \node at (-6.7,1.0) {\textcolor{purple!80!black}{Predictor}};
        
        % Draw a connecting line or arrow
        %\draw [->, dashed, purple!80!black] (pred.south) -- ++(0,-0.8);

        \node (plan1) at (pred1.east) [xshift=15pt] {\textcolor{purple!80!black}{\footnotesize \makecell{$k=2$}}};

        % Add this block for the second Speculation Depth Predictor
        \node (pred2) at (-1.98,1.0) [rectangle, draw=lightpurple!70, fill=lightpurple!70, minimum width=10pt, minimum height=10pt, inner sep=0pt, anchor=west] {};

        \node (plan2) at (pred2.east) [xshift=15pt] {\textcolor{purple!80!black}{\footnotesize \makecell{$k=3$}}};
        
        % approximation code
        \node (step1) at (-5,0) [rectangle, draw=lightblue!50, fill=lightblue!50, minimum width=25pt, minimum height=12pt, inner sep=0pt, anchor=west] {};
    
        \node (step1) at (-5,0) [rectangle, draw=lightblue!50, fill=lightblue!50, minimum width=25pt, minimum height=12pt, inner sep=0pt, anchor=west] {};
        \node (exc1) at (step1.east) [rectangle, draw=black!30, fill=black!30, minimum width=5pt, minimum height=12pt, inner sep=0pt, anchor=west] {};
        \node (step2) at (exc1.east) [rectangle, draw=lightblue!50, fill=lightblue!50, minimum width=20pt, minimum height=12pt, inner sep=0pt, anchor=west] {};
        \node (exc2) at (step2.east) [rectangle, draw=black!30, fill=black!30, minimum width=5pt, minimum height=12pt, inner sep=0pt, anchor=west] {};
        \node at (-6.3,0) {\textcolor{blue!80!black}{Approximation}};

        \node (plan1) at (step1.east) [yshift=15pt] {\textcolor{blue!80!black}{\scriptsize \makecell{A\\$\downarrow$}}};
        \node (plan2) at (step2.east) [yshift=15pt] {\textcolor{blue!80!black}{\scriptsize \makecell{B\\$\downarrow$}}};

        % target code
        \node (tar_step1) at ([yshift=-30pt]step1.west) [rectangle, draw=lightorange!50, fill=lightorange!50, minimum width=58pt, minimum height=12pt, inner sep=0pt, anchor=west] {};
        \node (tar_step2) at ([yshift=-45pt]step2.west) [rectangle, draw=lightorange!50, fill=lightorange!50, minimum width=55pt, minimum height=12pt, inner sep=0pt, anchor=west] {};
        \node at (-6.9,-1.6) {\textcolor{orange!80!black}{Target}};

        \draw [-, dashed, thick] ([yshift=7pt]pred1.west) -- node [yshift=-59.5pt] {\footnotesize $s_1$} ++(0,-3.75);
        \draw [-, dashed, thick] ([yshift=10pt]step2.west) -- node [yshift=-46pt] {\footnotesize $s_2$} ++(0,-2.9);
        \draw [-, dashed, thick] ([yshift=7pt]pred2.west) -- node [yshift=-59pt] {\footnotesize $s_3$} ++(0,-3.75);

        \node (tar_plan1) at (tar_step1.east) [yshift=15pt] {\textcolor{orange!80!black}{\scriptsize \makecell{A \textcolor{green!80!black}{\checkmark}\\$\downarrow$}}};

        % approximation code for episode 2
        \node (step3) at (tar_step2.east) [yshift=45pt] [rectangle, draw=lightblue!50, fill=lightblue!50, minimum width=25pt, minimum height=12pt, inner sep=0pt, anchor=west] {};
        \node (exc3) at (step3.east) [rectangle, draw=black!30, fill=black!30, minimum width=5pt, minimum height=12pt, inner sep=0pt, anchor=west] {};
        \node (step4) at (exc3.east) [rectangle, draw=lightblue!50, pattern=north west lines, pattern color=lightblue!100, minimum width=30pt, minimum height=12pt, inner sep=0pt, anchor=west] {};
        \node (exc4) at (step4.east) [rectangle, draw=black!30, pattern=north west lines, pattern color=black!50, minimum width=5pt, minimum height=12pt, inner sep=0pt, anchor=west] {};
        \node (step5) at (exc4.east) [rectangle, draw=lightblue!50, pattern=north west lines, pattern color=lightblue!100, minimum width=11pt, minimum height=12pt, inner sep=0pt, anchor=west] {};
        
        \node (plan3) at (step3.east) [yshift=15pt] {\textcolor{blue!80!black}{\scriptsize \makecell{C'\\$\downarrow$}}};
        \node (plan4) at (step4.east) [yshift=15pt] {\textcolor{blue!80!black}{\scriptsize \makecell{D'\\$\downarrow$}}};
        % \node (plan5) at (step5.east) [yshift=15pt] {\textcolor{blue!80!black}{\scriptsize \makecell{E\\$\downarrow$}}};

        % target code
        \node (tar_step3) at ([yshift=-30pt]step3.west) [rectangle, draw=lightorange!50, fill=lightorange!50, minimum width=78pt, minimum height=12pt, inner sep=0pt, anchor=west] {};
        \node (tar_exc3) at (tar_step3.east) [rectangle, draw=black!30, fill=black!30, minimum width=5pt, minimum height=12pt, inner sep=0pt, anchor=west] {};
        \node (tar_plan2) at (tar_step2.east) [yshift=15pt, xshift=-10pt] {\textcolor{orange!80!black}{\scriptsize \makecell{B \textcolor{green!80!black}{\checkmark}\\$\downarrow$}}};
        \node (tar_step4) at ([yshift=-45pt]step4.west) [rectangle, draw=lightorange!50, pattern=north west lines, pattern color=lightorange!100, minimum width=48pt, minimum height=12pt, inner sep=0pt, anchor=west] {};
        \node (tar_step5) at ([yshift=-60pt]step5.west) [rectangle, draw=lightorange!50, pattern=north west lines, pattern color=lightorange!100, minimum width=11pt, minimum height=12pt, inner sep=0pt, anchor=west] {};

        \draw [->, thick] (-7,-2.5) -- node [below, xshift=160pt, yshift=-3pt] {Time} (5.5,-2.5);
        \draw [-, dashed, thick] ([yshift=8pt]step4.west) -- node [yshift=-46pt] {} ++(0,-2.75);
        \draw [-, dashed, thick] ([yshift=8pt]step5.west) -- node [yshift=-46pt] {} ++(0,-2.77);

        \node (tar_plan3) at (tar_step3.east) [yshift=15pt, xshift=-22pt] {\textcolor{orange!80!black}{\scriptsize \makecell{C \textcolor{red}{\ding{55}}\\$\downarrow$}}};

        % online learning part
        \node (online) at ([yshift=75pt]tar_exc3.east) [rectangle, draw=almond!50, fill=almond!50, minimum width=78pt, minimum height=12pt, inner sep=0pt, anchor=west] {};
        \node (online_action) at (online.east) [yshift=18pt, xshift=-30pt] {\textcolor{orange!80!black}{\scriptsize \makecell{New data for Training:}}};
        \node (online_data) at (online.east) [yshift=10pt, xshift=27pt] {\textcolor{orange!80!black}{\scriptsize \makecell{($s = s_1, k = 3$)\\($s = s_2, k = 2$)\\($s = s_3, k = 1$)}}};

        \node at (-6.85,1.65) {\textcolor{almond!80!black}{Trainer}};

        % k = 3
        \node (pred3) at (tar_exc3.east) [yshift=58pt, rectangle, draw=lightpurple!70, fill=lightpurple!70, minimum width=10pt, minimum height=10pt, inner sep=0pt, anchor=west] {};
        \node (plan3) at (pred3.east) [xshift=15pt] {\textcolor{purple!80!black}{\footnotesize \makecell{$k=2$}}};

        % approximation code for episode 3
        \node (step6) at (tar_exc3.east)  [yshift=30pt, rectangle, draw=lightblue!50, fill=lightblue!50, minimum width=25pt, minimum height=12pt, inner sep=0pt, anchor=west] {};
        % cancel words
        \draw [-, solid, thick, draw=red, line width=1.5pt] ([yshift=6pt]step5.east) -- node [yshift=-55pt] {\footnotesize \textcolor{red}{\makecell{Cancel ongoing threads}}} ++(0,-2.7);
        % restart episode 3
        \node (exc6) at (step6.east) [rectangle, draw=black!30, fill=black!30, minimum width=5pt, minimum height=12pt, inner sep=0pt, anchor=west] {};
        \node (step7) at (exc6.east) [rectangle, draw=lightblue!50, fill=lightblue!50, minimum width=30pt, minimum height=12pt, inner sep=0pt, anchor=west] {};
        \node (exc7) at (step7.east) [rectangle, draw=black!30, fill=black!30, minimum width=5pt, minimum height=12pt, inner sep=0pt, anchor=west] {};
        

        \node (plan6) at (step6.east) [yshift=15pt] {\textcolor{blue!80!black}{\scriptsize \makecell{D\\$\downarrow$}}};
        \node (plan7) at (step7.east) [yshift=15pt] {\textcolor{blue!80!black}{\scriptsize \makecell{E\\$\downarrow$}}};

        % target code for episode 3
        \node (tar_step6) at (tar_exc3.east) [rectangle, draw=lightorange!50, fill=lightorange!50, minimum width=75pt, minimum height=12pt, inner sep=0pt, anchor=west] {};
        \node (tar_step7) at (exc6.east) [yshift=-45pt, rectangle, draw=lightorange!50, fill=lightorange!50, minimum width=63pt, minimum height=12pt, inner sep=0pt, anchor=west] {};

        \node (tar_plan6) at (tar_step6.east) [yshift=15pt] {\textcolor{orange!80!black}{\scriptsize \makecell{D \textcolor{green!80!black}{\checkmark}\\$\downarrow$}}};
        \node (tar_plan7) at (tar_step7.east) [yshift=15pt] {\textcolor{orange!80!black}{\scriptsize \makecell{E \textcolor{green!80!black}{\checkmark}\\$\downarrow$}}};

        \draw [-, dashed, thick] ([yshift=8pt]step5.west) -- node [yshift=-46pt] {} ++(0,-2.77);

        \draw [-, dashed, thick] ([yshift=8pt]step7.west) -- node [yshift=-45pt] {\footnotesize $s_5$} ++(0,-2.8);

        \draw [<->, thick]([yshift=5pt]step1.west |- 0,-3.5) -- node [below] {Total task time}   ([yshift=5pt]tar_step7.east |- 0,-3.5);

        \draw [-, dashed, thick] ([yshift=9pt]online.west) -- node [yshift=-69pt] {\footnotesize $s_4$} ++(0,-4.4);


        % curly bracket
        \draw[decoration={brace,raise=75pt},decorate] ([xshift=-2pt]step1.west) -- node[above=75pt] {episode 1} ([xshift=-5pt]step5.east);

        % \draw[decoration={brace,raise=105pt},decorate] (tar_step3.west) -- node[above=105pt] {episode 2} ([xshift=-5pt] tar_exc3.east);

        \draw[decoration={brace,raise=30pt},decorate] (online.west) -- node[above=30pt] {episode 2} ([yshift=90pt, xshift=2pt]tar_step7.east);
        
    \end{tikzpicture}
    }
    \caption{Dynamic Speculative Planning with Online Reinforcement Learning. The diagram illustrates how DSP adaptively adjusts speculation steps across planning episodes. (1) \textbf{Predictor} dynamically infers step values $k$ from current state $s_i$, while the \textbf{Approximation $(\approxi)$} and \textbf{Target $(\target)$} agents execute parallelly. (2) When mismatches occur (denoted by \textcolor{red}{\ding{55}}), ongoing threads are canceled and execution resumes from the last verified step. (3) An asynchronous \textbf{Trainer} collects state-step ($s, k$) pairs from completed episodes, continuously updating the predictor without blocking execution.}
    \label{fig:main_graph}
    % \vspace{-20pt}
    \vspace{-4.5pt}
    
    % \setlength{\belowcaptionskip}{-4pt} 
\end{figure}

% \caption{Dynamic Speculative Planing Pipeline}
% \label{fig:main_graph}
% \end{figure}

